\chapter{Introducción}

% --- SECCIÓN ------------------------------------------------------------------
% --- DEFINICIONES -------------------------------------------------------------
\section{Definiciones}

Denotaremos la n-ésima composición de una función $f$ consigo misma en un punto
$z$ con~$f^n(z)$.

\begin{definicion}[Órbitas]
  Dada una función $f$ y un punto $z$, la \emph{órbita hacia adelante} de $z$
  es el conjunto de puntos $z$, $f(z)$, $f^2(z)$, $f^3(z)$, $\ldots$ y la
  denotaremos por $\orbit^+(z)$. Si $f$ es un homeomorfismo, la \emph{órbita
  completa} de $z$, $\orbit(z)$, es el conjunto de puntos $f^n(z)$ con $n \in
  \Z$, y la \emph{órbita hacia} atrás de $z$, $\orbit^-(z)$, es el conjunto de
  puntos $z$, $f^{-1}(z)$, $f^{-2}(z)$,~\dots
\end{definicion}

\begin{definicion}[Punto fijo y periódico]
  Dada una función $f$ y un punto $\zeta$, $\zeta$ es un \emph{punto fijo} para
  $f$ si $f(\zeta) = \zeta$. El punto $\zeta$ es un \emph{punto periódico} de
  periodo $n$ si $f^n(\zeta) = \zeta$. El menor $n$ positivo tal que
  $f^n(\zeta) = \zeta$ se llama \emph{periodo primo} de $\zeta$. Denotaremos el
  conjunto de puntos periódicos de periodo $n$ (no necesariamente primo) por
  $\periodicpoints{n}(f)$, y el conjunto de puntos fijos por $\fixedpoints(f)$.
  El conjunto de todas las iteraciones de un punto periódico es una
  \emph{órbita periódica}.
\end{definicion}

\begin{definicion}[Punto eventualmente periódico]
  Un punto $z$ es \emph{eventualmente periódico} de periodo $n$ si $z$ no es
  periódico pero existe $m > 0$ tal que $f^{i+n}(z) = f^i(z)$ para todo $i \geq
  m$. Es decir, $f^i(z)$ es periódico si $i \geq m$.
\end{definicion}

\begin{definicion}[Punto asintótico y conjunto estable]
  Sea $p$ periódico de periodo $n$. Un punto $z$ es \emph{asintótico hacia
  adelante} a $p$ si $\lim_{i \to \infty} f^{in}(z) = p$. El \emph{conjunto
  estable} de $p$, denotado por $W^s(p)$, está formado por todos los puntos
  asintóticos hacia adelante a~$p$.
\end{definicion}

Si $p$ es no periódico, también podríamos definir puntos asintóticos hacia
adelante de la siguiente forma, $|f^i(z) - f^i(p)| \to 0|$ cuando $i \to
\infty$. Además, si $f$ tiene inversa, podemos considerar puntos
\emph{asintóticos hacia atrás} haciendo tender $i \to -\infty$ en la definición
de encima. Al conjunto de puntos asintóticos hacia atrás a $p$ se le llama el
conjunto \emph{inestable} de $p$ y se denota por~$W^u(p)$.

\begin{definicion}[Punto crítico]
  Un punto $z$ es un \emph{punto crítico} de $f$ si $f'(z) = 0$. El punto
  crítico es no degenerado si $f''(z) \neq 0$. El punto crítico es degenerado
  si~$f''(z) = 0$.
\end{definicion}

\begin{definicion}[Función Racional]
  Una \emph{función racional} es una función $R: \ExtC \to \ExtC$ de la forma
  \begin{equation*}
    R(z) = \frac{P(z)}{Q(z)},
  \end{equation*}
  donde $P(z)$ y $Q(z)$ son polinomios, sin ser ambos nulos.
\end{definicion}

Supongamos que ni $P$ ni $Q$ son el polinomio nulo, en este caso $R$ no está
definida en las raíces comunes de $P$ y $Q$. Sin embargo, si estas existen,
podemos cancelar los factores comunes de $P$ y $Q$, obteniendo así dos
polinomios que no tendrán raíces comunes $P_1$ y $Q_1$. De ahora en adelante
consideraremos que $P$ y $Q$ no tienen raíces comunes, y si las tuvieran,
consideraríamos $R = P_1 / Q_1$. En estas circunstancias seguiremos los
convenios habituales con el punto del infinito en la esfera de Riemann: si $Q$
es el polinomio nulo, entonces $R$ es la función constante $\infty$; si $Q(z) =
0$ y $P$ no es el polinomio nulo, entonces $R(z) = \infty$; y tomaremos
$R(\infty) = \lim_{z \to \infty} R(z)$.

\begin{definicion}[Métrica cordal]
  En $\ExtC$ definimos la métrica cordal de la siguiente manera:
  \begin{equation*}
    \sigma(z, w) = \frac{|z - w|}{\sqrt{1 + |z|^2}\sqrt{1 + |w|^2}}, \quad \sigma(z, \infty) = \frac{1}{\sqrt{1 + |z|^2}}
  \end{equation*}
\end{definicion}

% --- SECCIÓN ------------------------------------------------------------------
% --- TRANSFORMACIONES DE MÖBIUS -----------------------------------------------
\section{Transformaciones de Möbius}

\begin{definicion}[Transformación de Möbius]
  Una transformación de Möbius es una función racional de la forma
  \begin{equation*}
    R(z) = \frac{az + b}{cz + d}, \quad ad-bc \neq 0,
  \end{equation*}
  donde $z$, $a$, $b$, $c$ y $d$ son números complejos.
\end{definicion}

Siguiendo los convenios habituales con el punto del infinito en la esfera de
Riemann, si $c \neq 0$: $R(\infty) = a/c$ y $R(-d/c) = \infty$. Si $c = 0$,
entonces $R(\infty) = \infty$.

\begin{definicion}[Funciones Racionales Conjugadas]
  Sean $R$ y $S$ dos funciones racionales. Decimos que $R$ y$S$ son conjugadas
  si existe alguna transformación de Möbius $g$ tal que
  \begin{equation*}
    S = g R g^{-1}.
  \end{equation*}
\end{definicion}

La propiedad de ser conjugadas es una relación de equivalencia. Además, dos
funciones racionales conjugadas $R$ y $S$ cumplen las siguientes propiedades:
\begin{enumerate}[label=(\arabic*)]
  \item ambas funciones tienen el mismo grado;
  \item $S^n = g R^n g^{-1}$; y,
  \item un punto $\zeta$ es punto fijo de $R$ si y solo si $g(\zeta)$ lo es
        de~$S$.
\end{enumerate}

\begin{proposicion}
  Sea $R$ una transformación de Möbius distinta de la identidad. Entonces, $R$
  tiene dos puntos fijos en $\ExtC$, contando multiplicidades.
\end{proposicion}
\begin{proof}
  Un punto $\zeta$ de $\ExtC$ es un punto fijo de $R$ si $R(\zeta) = \zeta$.
  Distinguimos dos casos:
  \begin{enumerate}[label=(\roman*)]
    \item si $c$ es cero y $a \neq d$, los puntos fijos son infinito y la
          solución de $(d-a)\zeta = b$; \label{item:mobius_inf_fijo}
    \item si $c$ es cero y $a = d$, el infinito es un punto fijo doble; y,
          \label{item:mobius_inf_fijo_doble}
    \item si $c$ es no nulo, los puntos fijos son las dos soluciones (contando
          multiplicidades) de $c\zeta^2 + (d - a)\zeta - b = 0$.
  \end{enumerate}
\end{proof}

\begin{proposicion}
  Sea $R$ una transformación de Möbius distinta de la identidad. Entonces, $R$
  es conjugada a una función de una de las siguientes formas:
  \begin{enumerate}[label=(\arabic*)]
    \item $S(w) = w + 1$, si $R$ tiene un punto fijo con multiplicidad dos; o
    \item $S(w) = qw$, donde $w \in \C - \{0, 1\}$, si $R$ tiene dos puntos
          fijos diferentes.
  \end{enumerate}
\end{proposicion}
\begin{proof}
  Supongamos $R$ tiene un único punto fijo $\zeta$ con multiplicidad dos.
  Primero, si $\zeta$ es finito, lo enviamos al infinito mediante la
  transformación de Möbius $g(z) = 1/(z - \zeta)$. La conjugada $S = g R
  g^{-1}$, que también es una transformación de Möbius, tiene como único punto
  fijo al infinito. Por tanto, es de la forma $w \mapsto aw + b$, con $a = 1$
  por \ref{item:mobius_inf_fijo_doble} de la demostración anterior, y $b \neq
  0$ (porque $R$ no es la identidad). Así $S(w) = w + b$, y conjugando con
  $h(w) = w/b$ tenemos que $R$ es conjugada a $w \mapsto w + 1$.

  Demostremos ahora el caso en que $R$ tiene dos puntos fijos diferentes,
  $\zeta_1$ y $\zeta_2$. En este caso enviaremos un punto fijo al cero y el
  otro al infinito. Si ambos puntos son finitos, utilizamos la transformación
  $g(z) = (z - \zeta_1)/(z - \zeta_2)$. Si uno de los dos ya es infinito,
  digamos $\zeta_2$, utilizaremos la transformación $g(z) = z - \zeta_1$. La
  transformación conjugada $S = g R g^{-1}$ fija el cero y el infinito, luego
  es de la forma $w \mapsto aw + b$, con $a \in \C - \{0, 1\}$ y $b = 0$ por
  \ref{item:mobius_inf_fijo} de la demostración anterior. Concluyendo que $R$ es
  conjugada a $S = aw$, con $a \in \C - \{0, 1\}$.
\end{proof}
